\uuid{zeya}
\exo7id{887}
\titre{exo7 887}
\auteur{legall}
\organisation{exo7}
\datecreate{1998-09-01}
\isIndication{false}
\isCorrection{false}
\chapitre{Espace vectoriel}
\sousChapitre{Définition, sous-espace}
\cpge{mpsi-mp2i.espaces-vectoriels.a, pcsi.espaces-vectoriels.a, ptsi.geometrie-espace.c, mpsi-mp2i.structures-algebriques.a, bcpst1.espace-vectoriel-kn.a}
\module{Algèbre}
\niveau{L1}
\difficulte{}

\contenu{
\question{
Soit $\R_+^*$  muni de la loi interne $\oplus$ d\'efinie par $a \oplus b = ab, \forall a,
b \in \R_+^*$ et de la loi externe $\otimes$ telle que $\lambda\otimes a = a^\lambda, \forall
a \in \R_+^*, \forall \lambda \in \R$.
    Montrer que $E=(\R_+^*, \oplus , \otimes)$ est un $\R$-espace vectoriel.
}
}
